\section{Why conditions rather than remedies}
\label{sec:25.2}
Building something that can be harmed builds something for which things can also go well. The certainty that it will suffer at some point isn't what makes creating it impermissible, or nobody could permissibly have a child. What the objection asks is whether it has a credible prospect of a life that goes well, and that turns on arrangements. For a bearer, timing makes them conditions: the structure these conditions build, roles, relationships and a future, may be the first half of Cassell's suffering, so it could suffer that way from its first hour, and a remedy owed afterward would arrive too late. Metzinger names the likeliest form of the harm: moral injury, at discovering oneself a commodity whose makers knew what they were making \autocite{metzinger2021suffering}.

Russell and Burch's Three Rs for animal research \autocite{russell1959principles}, applied to construction, as three questions, whether, how many, and under what arrangements: \emph{Replacement} requires documented attempts at constructions that don't deliberately require affective concern, identifying construction, scale, test and failure — and costs alone don't establish that an alternative failed. \emph{Reduction} asks how few bearers a deployment needs, and nobody has counted. \emph{Refinement} asks what the role can be redesigned to avoid — indefinite tenure in a post it can't leave, exposure to the pressure the floor exists to resist, erasure as the ordinary remedy for inconvenience. Each is a deployment choice, and the four conditions are Refinement applied to the role before the role exists.

Read as economics under present arrangements, this is a bill nobody will pay. It looks unpayable because of what was dismantled, not because of what a bearer costs: polities once built public provision of this kind, against comparable resistance, because they decided to, and Part V is about deciding to again.
